\section{Supervised Fine-Tuning}
\label{sec:sft}

\subsection{Model Selection}
\label{sec:sft:model}

We evaluated two open-source instruction-tuned models as base candidates: Llama~3.1~8B~\cite{roziere2024codellama} and Qwen3-8B~\cite{qwen2025qwen3}, both in 4-bit quantized form via Unsloth. Qwen3-8B achieves a substantially higher zero-shot baseline on our reliability engineering questions (70.7\% vs.\ 37.5\%), likely due to its stronger mathematical pre-training. Moreover, fine-tuning Llama~3.1~8B caused a regression on numeric questions ($-1.4$pp) while improving only format-dependent categories (formula $+28.6$pp, text $+15.8$pp), suggesting catastrophic forgetting of mathematical capabilities. We therefore selected Qwen3-8B as our base model for all subsequent experiments.

By analyzing the model's thinking traces, we observed that the base model uses these blocks as \emph{scratch work}: it explores solution avenues, backtracks when an approach seems wrong, and iterates toward the answer, much like a human solving a problem for the first time. Our SFT training data, by contrast, consists of polished step-by-step solutions that proceed directly to the correct answer without false starts. Training on these clean reasoning chains with thinking enabled would effectively require the model to ``already know the answer'' and go straight to a proof, conflicting with the exploratory nature of the thinking trace. We found that this mismatch caused the model to overfit, producing empty thinking tags followed by shallow pattern-matched answers. Disabling thinking during both training and inference resolved this: the model generates direct chain-of-thought responses consistent with the training format, while preserving its internal reasoning capabilities for use during GRPO (Section~\ref{sec:grpo}).

\subsection{Training Setup}
\label{sec:sft:setup}

\subsubsection{LoRA Configuration}

We use 4-bit QLoRA~\cite{dettmers2023qlora} via the Unsloth framework, which provides optimized CUDA kernels for approximately 40\% faster training. LoRA adapters~\cite{hu2022lora} are applied to all attention and MLP projection layers. Key hyperparameters are rank $r=16$, scaling factor $\alpha=32$, and dropout $0.05$.

\subsubsection{Training Hyperparameters}

Table~\ref{tab:training_config} summarizes the training configuration. NEFTune~\cite{jain2023neftune} adds uniform noise to embedding vectors during training, serving as a regularizer against memorization on small datasets. We found $\alpha=5$ to be optimal: $\alpha=7$ degraded performance by $-1.3$pp, and $\alpha=10$ showed only marginal gains. Early stopping with patience 2 (monitoring validation loss on a 10\% held-out split) consistently converged to epoch 4 regardless of the maximum epoch setting (6 or 8), indicating a stable training dynamic.

\begin{table}[t]
\centering
\small
\caption{SFT training configuration (best setup).}
\label{tab:training_config}
\begin{tabular}{@{}ll@{}}
\toprule
Parameter & Value \\
\midrule
Learning rate & $2 \times 10^{-4}$ (cosine) \\
Optimizer & AdamW 8-bit \\
Batch size & 2 (grad.\ accum.\ 4, eff.\ 8) \\
Max epochs & 8 (early stop $\to$ epoch 4) \\
Warmup steps & 5 \\
Weight decay & 0.01 \\
NEFTune $\alpha$ & 5 \\
Precision & BF16 \\
Max sequence length & 2048 tokens \\
\bottomrule
\end{tabular}
\end{table}

\subsubsection{Data Format}

Training examples follow a three-turn chat format: a \emph{system} message establishing the reliability engineering expert persona, a \emph{user} message containing the question, and an \emph{assistant} message containing the full chain-of-thought reasoning followed by the final answer. The tokenizer's \texttt{apply\_chat\_template()} handles model-specific formatting.

\subsection{Evaluation Protocol}
\label{sec:sft:eval}

\subsubsection{K-Fold Cross-Validation}

Given the small dataset sizes (215--866 examples), we use $k$-fold cross-validation rather than a single train/test split to obtain reliable performance estimates. Initial experiments used 5-fold CV (172 train / 43 test per fold), but 5 folds limit the minimum achievable $p$-value to 0.0625 under the Wilcoxon test, precluding statistical significance at $p < 0.05$. We therefore moved to 10-fold CV for experiments on 866 samples ($\sim$780 train / $\sim$87 test per fold), which achieves $p = 0.002$ when all folds show improvement.

For each fold, a fresh base model is loaded and fine-tuned from scratch to avoid information leakage between folds. All evaluation uses greedy decoding (\texttt{do\_sample=False}) with \texttt{max\_new\_tokens=4096}, ensuring deterministic and reproducible results. Early experiments with stochastic decoding inflated improvements (e.g., $+7.0\%$ stochastic vs.\ $+2.3\%$ greedy).

\subsubsection{Answer Comparison}

We extract the final answer from model responses using pattern matching (``Final Answer:'', ``The answer is:'', ``Therefore,'') with fallback to the last non-empty line. Extracted answers are then compared against ground truth through a multi-level pipeline:

\begin{enumerate}
    \item \textbf{Exact match:} Normalized strings (lowercase, strip whitespace, remove LaTeX wrappers such as \verb|\boxed{}|, convert \verb|\frac{a}{b}| to \verb|a/b|).
    \item \textbf{Numerical match:} Extract all numbers via regex; accept if all values match within 5\% relative tolerance.
    \item \textbf{Fraction-to-decimal:} Convert fractions (e.g., 14/33 $\leftrightarrow$ 0.4242).
    \item \textbf{Percentage conversion:} Normalize percentages (e.g., 25.5\% $\leftrightarrow$ 0.255).
\end{enumerate}

An answer is considered correct if any of the above criteria is satisfied.

\subsubsection{Statistical Testing}

We use the Wilcoxon signed-rank test to compare per-fold baseline and fine-tuned accuracies. This non-parametric paired test is appropriate for small paired samples where normality cannot be assumed. We use $p < 0.05$ as the significance threshold. We additionally track \emph{question flips}---questions that change correctness between baseline and fine-tuned models---as a measure of improvement robustness.

\subsection{SFT Results}
\label{sec:sft:results}

Table~\ref{tab:sft_results} presents the main SFT results across dataset sizes. The improvement scales non-linearly with dataset size: from $+2.3\%$ (not significant) on 215 original questions to $+18.6\%$ ($p=0.002$) on 866 augmented questions. All 10 folds showed improvement, with per-fold deltas ranging from $+12.6\%$ to $+26.4\%$. This suggests a threshold effect around 500 examples, below which LoRA fine-tuning struggles to overcome the regularization of frozen base weights.

\begin{table}[t]
\centering
\small
\caption{SFT results across dataset sizes (best config per size). All use Qwen3-8B, greedy decoding, LR=$2\!\times\!10^{-4}$, NEFTune=5.}
\label{tab:sft_results}
\begin{tabular}{@{}rcccccc@{}}
\toprule
$N$ & Folds & Ep. & Base & FT & $\Delta$ & $p$ \\
\midrule
215 & 5 & 4          & 70.7 & 73.0 & $+2.3$  & 0.50 \\
280 & 5 & 4          & 71.0 & 73.0 & $+2.0$  & 0.69 \\
501 & 5 & 4          & 59.1 & 65.3 & $+6.2$  & 0.063 \\
\textbf{866} & \textbf{10} & \textbf{8$\to$4} & \textbf{63.6} & \textbf{82.2} & $\boldsymbol{+18.6}$ & \textbf{0.002} \\
\bottomrule
\end{tabular}
\end{table}

The question flip ratio is particularly informative: at 215 samples, the model trades nearly as many correct answers as it gains (1.3:1, 23~W$\to$R vs.\ 18~R$\to$W), while at 866 samples, gains overwhelm regressions (4.6:1, 206~W$\to$R vs.\ 45~R$\to$W), confirming that paraphrase augmentation produces robust improvements. Note that baseline accuracy is lower on larger datasets (63.6\% on 866 vs.\ 70.7\% on 215) because augmented test sets include paraphrased questions the base model has not seen.

\subsection{Ablation Studies}
\label{sec:sft:ablation}

Table~\ref{tab:ablations} summarizes key ablation results on the 215-question dataset (5-fold CV, greedy decoding).

\begin{table}[t]
\centering
\small
\caption{SFT ablations on 215 questions (5-fold CV, greedy).}
\label{tab:ablations}
\begin{tabular}{@{}lcc@{}}
\toprule
Variation & $\Delta$ & $p$ \\
\midrule
\textbf{Best (LR=$\mathbf{2\!\times\!10^{-4}}$, NEF=5, 4ep)} & $\boldsymbol{+2.3}$ & 0.50 \\
Lower LR ($1\!\times\!10^{-4}$) & $-0.5$ & 1.0 \\
Higher NEFTune ($\alpha\!=\!7$) & $-1.3$ & 0.38 \\
Lower LoRA rank ($r\!=\!8$, $\alpha\!=\!16$) & $+1.4$ & 0.88 \\
3 epochs (vs.\ 4) & $+0.9$ & 0.88 \\
5 epochs (vs.\ 4) & $+1.9$ & 0.50 \\
DPO (instead of SFT) & $-1.7$ & 0.63 \\
SFT + GRPO (early) & $+0.5$ & 0.88 \\
\bottomrule
\end{tabular}
\end{table}

LR=$2\!\times\!10^{-4}$ consistently outperforms $1\!\times\!10^{-4}$, consistent with the recommendation that LoRA adapters benefit from learning rates approximately 10$\times$ higher than full fine-tuning. Early stopping consistently selects epoch 4 across all dataset sizes. Both DPO ($-1.7\%$) and early GRPO ($+0.5\%$) underperform pure SFT on 215 questions; the early GRPO used TRL with LLM-judge rewards, motivating our improved pipeline with verifiable rewards on a larger training set (Section~\ref{sec:grpo}).
